\subsection{Imagen directa de un $\mathcal{A}$-Módulo}
\label{subsection:0.4.2-fr}

\begin{env}[4.2.1]
\label{0.4.2.1-fr}
Sean $(X,\mathcal{A})$, $(Y,\mathcal{B})$ dos espacios anillados y
$\Psi=(\psi,\theta)$ un morfismo
$(X,\mathcal{A})\to(Y,\mathcal{B})$. Entonces $\psi_*(\mathcal{A})$ es un
haz de anillos sobre $Y$ y $\theta$ un homomorfismo
$\mathcal{B}\to\psi_*(\mathcal{A})$ de haces de anillos. Sea ahora
$\mathcal{F}$ un $\mathcal{A}$-Módulo. Su imagen directa
$\psi_*(\mathcal{F})$ es un haz de grupos abelianos sobre $Y$. Además,
para todo abierto $U\subset Y$,
\[
  \Gamma(U,\psi_*(\mathcal{F}))
  =\Gamma(\psi^{-1}(U),\mathcal{F})
\]
está provisto de una estructura de módulo sobre el anillo
$\Gamma(U,\psi_*(\mathcal{A}))=\Gamma(\psi^{-1}(U),\mathcal{A})$.
Como las aplicaciones bilineales que definen estas estructuras son
compatibles con las operaciones de restricción, definen sobre
$\psi_*(\mathcal{F})$ una estructura de $\psi_*(\mathcal{A})$-Módulo.
El homomorfismo $\theta:\mathcal{B}\to\psi_*(\mathcal{A})$ permite entonces
definir también sobre $\psi_*(\mathcal{F})$ una estructura de
$\mathcal{B}$-\emph{Módulo}. Diremos que este $\mathcal{B}$-Módulo es la
\emph{imagen directa de $\mathcal{F}$ por el morfismo $\Psi$}, y lo
denotaremos por $\Psi_*(\mathcal{F})$. Si $\mathcal{F}_1$, $\mathcal{F}_2$
son dos $\mathcal{A}$-Módulos sobre $X$ y $u$ es un
$\mathcal{A}$-homomorfismo $\mathcal{F}_1\to\mathcal{F}_2$, es inmediato
(al considerar las secciones sobre los abiertos de $Y$) que $\psi_*(u)$
es un $\psi_*(\mathcal{A})$-homomorfismo
$\psi_*(\mathcal{F}_1)\to\psi_*(\mathcal{F}_2)$ y, \emph{a fortiori}, un
$\mathcal{B}$-homomorfismo
$\Psi_*(\mathcal{F}_1)\to\Psi_*(\mathcal{F}_2)$. Como
$\mathcal{B}$-homomorfismo, lo denotaremos por $\Psi_*(u)$. Se ve, pues, que
$\Psi_*$ es un \emph{funtor covariante} de la categoría de los
$\mathcal{A}$-Módulos a la de los $\mathcal{B}$-Módulos. Además, es
inmediato que este funtor es \emph{exacto por la izquierda} (G, II, 2.12).

En $\psi_*(\mathcal{A})$, la estructura de $\mathcal{B}$-Módulo y la
estructura de haz de anillos definen una estructura de
$\mathcal{B}$-Álgebra. Se denotará por $\Psi_*(\mathcal{A})$ esta
$\mathcal{B}$-Álgebra.
\end{env}
